\subsection{The exact missing term in the historical covariance-score criterion}
\label{sec:v17-mean-shift-decomposition}
The distinction between a validation-class-centred Gaussian
covariance score and a genuine source-trained-mean predictive
log score can be made \emph{exact}, rather than relying only
on a synthetic translation counterexample.

For one held \emph{source} validation person $g$ with $n_g$
examples in $p$ dimensions, let
$\widehat\mu_{g,c}$ denote its empirical class means,
$\widehat\mu_{\mathrm{train},c}$ the means fitted exclusively
on the source-training people, and
$\Delta_{g,c}=\widehat\mu_{g,c}-
\widehat\mu_{\mathrm{train},c}$. For any common,
source-fitted positive definite regularised covariance
$A_\lambda$, the two V16 criteria satisfy
\begin{equation}
\boxed{
 C_{\mathrm{pred},g}(\lambda)
 =
 C_{\mathrm{cent},g}(\lambda)
 +\frac{1}{p}\sum_{c\in\{0,1\}}
 \frac{n_{g,c}}{n_g}
 \Delta_{g,c}^\top A_\lambda^{-1}\Delta_{g,c}.
}
\label{eq:v17-exact-nll-shift}
\end{equation}
The additional term is \emph{nonnegative} and
becomes zero exactly when the validation class means
match the training class means, for all classes present.
It is a covariance-weighted \emph{source-person
class-mean shift penalty}, not a regularisation
term chosen by target data.

\emph{Proof.} For each validation trial of class $c$,
decompose
$x_{gi}-\widehat\mu_{\mathrm{train},c}=
(x_{gi}-\widehat\mu_{g,c})+\Delta_{g,c}$.
Expand the quadratic form under $A_\lambda^{-1}$.
The cross term vanishes within each class because
$\sum_{i:y_i=c}(x_{gi}-\widehat\mu_{g,c})=0$.
Both log-determinant terms agree, and the remaining
class-mean quadratic terms give
Eq.~\eqref{eq:v17-exact-nll-shift}.\hfill$\square$

The penalty is $\lambda$-dependent because it measures
class-mean drift in the metric induced by $A_\lambda^{-1}$.
Its derivative can have \emph{either} sign:
\[
 \frac{\partial}{\partial\lambda}
 \bigl(C_{\mathrm{pred},g}-C_{\mathrm{cent},g}\bigr)
 =-\frac1p\sum_c\frac{n_{g,c}}{n_g}
 \Delta_{g,c}^\top A_\lambda^{-1}
 (\nu I-S)A_\lambda^{-1}\Delta_{g,c}.
\]
It follows that the proper predictive criterion
can select more \emph{or} less shrinkage than
the class-centred comparator. Merely asserting
that predictive selection is stronger or always
more regularised is unsupported. An independent
executable test verifies the identity for 300
randomly generated full-rank covariance operators,
multiple feature dimensions and source/validation
class-mean shifts, without using any held outer
target EEG data.

The identity provides a mechanistic interpretation
of the V16 estimand audit: the old frozen covariance
selection scores how well $A_\lambda$ describes
\emph{within-class residual fluctuations} in a
held source person; the correctly predictive version
also penalises \emph{between-source-person class-mean
misalignment}. This is not a newly invented
covariance estimator, and it does not establish that
minimising the added penalty improves held-person AUC.
A real, correctly nested evaluation of this second
source-selection criterion remains essential.
